\begin{appendices}

\chapter{Self Appraisal}

The project plan followed incremental delivery. Low complexity experiments were stabilised before new simulator systems were added.

Research discipline improved through controlled benchmarking. Experiments now use fixed seeds with matched budgets from the start.

\section{Critical self evaluation}
Scope remained tight when implementation risk increased. This decision improved result reliability.

Early writing pace can improve in future projects. Too much time was spent in exploration before full chapter drafting.

\section{Personal reflection and lessons learned}
Reproducibility emerged as a core engineering skill. Simple baselines also exposed design issues quickly.

\section{Legal social ethical and professional issues}
Public research sources were used with citation. No personal data was collected from participants.

Energy use was considered during long training runs. Waste was reduced by using short smoke checks before long runs.

\subsection{Legal issues}
Open publication sources were followed and cited clearly.

\subsection{Social issues}
This project did not interact with human participants.

\subsection{Ethical issues}
Conclusions were not overstated in limited scope experiments.

\subsection{Professional issues}
Version control was used with traceable experiment settings.

\chapter{External Material}
Publicly available papers were used as academic background material.

FastF1 based track data from public tooling was used in the workflow.

\section{Related Work Comparison}
Table~\ref{tab:related-work-comparison} summarises key studies reviewed in Chapter 1 and highlights how this dissertation positions its contribution.

\begin{table}[htbp]
\centering
\scriptsize
\caption{Comparison of related work and relevance to this dissertation}
\label{tab:related-work-comparison}
\resizebox{\textwidth}{!}{
\begin{tabular}{|p{2.7cm}|p{2.8cm}|p{2.7cm}|p{2.5cm}|p{2.5cm}|p{2.7cm}|p{2.8cm}|}
\hline
\textbf{Study} & \textbf{Domain and simulator fidelity} & \textbf{Decision method} & \textbf{Uncertainty treatment} & \textbf{Evaluation rigour} & \textbf{Limitations} & \textbf{Direct relevance to thesis} \\
\hline
Heilmeier et al. (2018) & Circuit motorsport with lap-wise race simulation and tyre and fuel effects & Strategy simulation and what-if analysis & Includes overtaking and degradation effects in model dynamics & Practical validation on historical race setting & Not focused on RL benchmark fairness across seeds & Motivates staged simulator design with tyre extension \\
\hline
Heilmeier et al. (2020) & Circuit motorsport strategy support with ANN components & Surrogate modelling and optimisation support & Captures race variability through modelled scenarios & Strong engineering framing and deployability discussion & Limited focus on RL reproducibility protocol & Supports need for interpretable strategy tools \\
\hline
Piccinotti et al. (2021) & Formula 1 strategy identification in planning setting & Online planning methods & Dynamic race context in planning loop & Demonstrates value of planning for strategy choice & Limited DQN-family comparative analysis & Clarifies boundary between planning and RL benchmarking goals \\
\hline
Liu and Fotouhi (2020) & Formula E race strategy with energy constraints & ANN plus MCTS & Multi-stage uncertainty in energy and race conditions & Demonstrates decision quality in time-critical domain & Not designed as controlled RL reproducibility benchmark & Supports tactical decision framing under uncertainty \\
\hline
Liu et al. (2024) & Formula E multi-car strategy learning context & Reinforcement learning based strategy development & Competitive race uncertainty in multi-car setting & Strong domain result emphasis & Less emphasis on seed-controlled statistical protocol & Provides adjacent motorsport RL evidence base \\
\hline
Mnih et al. (2015), van Hasselt et al. (2016), Wang et al. (2016), Hessel et al. (2018) & General RL benchmarks & DQN, Double DQN, Dueling DQN, Rainbow & Algorithmic robustness mechanisms in value learning & Benchmark-oriented algorithm studies & Not motorsport-specific and limited tactical interpretability & Defines core algorithm comparison set for this dissertation \\
\hline
Henderson et al. (2018), Agarwal et al. (2021) & Cross-domain deep RL evaluation & Reproducibility and statistical evaluation methodology & Focus on stochastic outcome variation and seed effects & Strong critique of weak reporting practice & Do not provide motorsport-specific protocol & Directly motivates fixed-seed and interval-based evaluation design \\
\hline
Sankary and Ostfeld (2018) & Robust optimisation in water systems & Evolutionary optimisation with stochastic scenario selection & Explicit stochastic scenario resampling & Robustness-focused evaluation method & Different optimisation paradigm from RL & Supports uncertainty-aware evaluation discipline \\
\hline
This dissertation & F1-inspired competitive simulator with staged complexity & Reproducible DQN-family benchmark with behavioural diagnostics & Controlled stochasticity pathway with tyre-related uncertainty stage & Matched budget and seed-controlled reporting with confidence intervals & Simplified simulator abstraction and staged feature scope & Targets interpretable and statistically grounded method comparison \\
\hline
\end{tabular}
}
\end{table}

\end{appendices}
